% Pregunta 314 (2019-febrero-OEP2018, n.º 25): fronteras de posibilidades de producción de A y B con 40 horas; son paralelas (1 cereal = 2 telas en los dos), B tiene ventaja absoluta pero no hay ventaja comparativa. Autor: Víctor Gutiérrez Marcos
\documentclass[tikz,border=6pt]{standalone}
\input{tikz-tcee}
\begin{document}
\begin{tikzpicture}[tcee,x=1cm,y=0.5cm]
  \ejes{Cereal}{Tela}{6.4}{11.6}
  % Fronteras con 40 horas: A (10 h por cereal, 5 h por tela), B (8 h y 4 h)
  \draw[curva1] (0,8) -- (4,0);
  \draw[curva2] (0,10) -- (5,0);
  % Cortes
  \node[below] at (4,0) {$4$};
  \node[below] at (5,0) {$5$};
  \node[left] at (0,8) {$8$};
  \node[left] at (0,10) {$10$};
  % Triángulos de pendiente: 1 cereal más = 2 telas menos en los dos
  \draw[guia,tceemorado] (0.8,6.4) -- (0.8,4.4) -- (1.8,4.4);
  \node[etiqueta,tceemorado,font=\footnotesize,below] at (1.3,4.4) {$1$};
  \node[etiqueta,tceemorado,font=\footnotesize,left] at (0.8,5.4) {$2$};
  \draw[guia,tceeazul] (3.4,3.2) -- (4.4,3.2) -- (4.4,1.2);
  \node[etiqueta,tceeazul,font=\footnotesize,above] at (3.9,3.2) {$1$};
  \node[etiqueta,tceeazul,font=\footnotesize,right] at (4.4,2.2) {$2$};
  % Rótulos de los precios relativos de autarquía
  \node[etiqueta,tceemorado,anchor=west,align=left] at (0.2,2.4) {$\mathrm{FPP}_A$\\$P_C/P_T=2$};
  \node[etiqueta,tceeazul,anchor=west,align=left] at (2.3,7.4) {$\mathrm{FPP}_B$\\$P_C/P_T=2$};
  \node[etiqueta,font=\footnotesize,align=left,anchor=west] at (2.9,10.6) {misma pendiente:\\sin ventaja comparativa};
\end{tikzpicture}
\end{document}
